\section{La santé mentale et la charge de travail des étudiants}\label{la-sante-mentale-et-la-charge-de-travail-des-etudiants}

\stanceinfo{\textbf{Congrès de la FCEG 2018} [2018-01-07]}{2024-01-02}

\officialstance{La Fédération canadienne étudiante de génie estime que les étudiants en ingénierie sont confrontés à des problèmes de santé mentale en raison de la structure et de la charge de travail de leurs programmes. Toutes les parties prenantes, y compris le corps enseignant, les organismes d'accréditation et les leaders étudiants, ont la responsabilité de revoir et de modifier leurs pratiques afin d'offrir aux étudiants un environnement d'apprentissage sain et efficace.}

\stanceheading{La position des étudiants}
\begin{itemize}
 \item Les étudiants en ingénierie présentent un risque élevé d'effets négatifs sur la santé mentale, avec quatre thèmes récurrents présents en littérature [3];
 \begin{itemize}
  \item La charge de travail des ingénieurs est un facteur de stress déterminant;
  \item Il existe des obstacles qui empêchent les étudiants en ingénierie de chercher de l'aide pour des problèmes de santé mentale;
  \item Les pairs s'appuient davantage les uns sur les autres pour faire face au stress et aux problèmes de santé mentale;
  \item La pandémie de COVID-19 a exacerbé les problèmes de santé mentale des étudiants.
 \end{itemize}
 \item Les étudiants en ingénierie se déclarent très stressés, anxieux et dépressifs, certains documents faisant état de 28 \% de stress modéré à extrêmement grave, de 35 \% d'anxiété modérée à extrêmement grave et de 34 \% de dépression modérée à extrêmement grave, d'après une enquête réalisée en 2021 [4].
 \item La perception généralisée selon laquelle un niveau de stress élevé et une mauvaise santé mentale sont normaux et prévisibles dans le domaine de l'ingénierie est préjudiciable et ne tient pas compte des problèmes de santé mentale des étudiants de premier cycle actuels et de ceux qui veulent choisir l'ingénierie.
 \item La charge de travail excessive des étudiants est probablement l'une des principales causes de ces effets négatifs sur la santé mentale, et des mesures devraient être prises pour réduire le stress universitaire lié à la charge de travail.
 \item Les problèmes de santé mentale touchent de manière disproportionnée les femmes et les autres minorités, notamment les personnes LGBTQ2S+, les autochtones et les personnes de couleur [4].
 \item Les étudiants en ingénierie sont peu enclins à rechercher un soutien en matière de santé mentale, même s'ils déclarent eux-mêmes des niveaux plus élevés de détresse en matière de santé mentale, en raison de différents facteurs, notamment physiques et culturels [5].
\end{itemize}

\stanceheading{Les recommandations aux partenaires, aux parties intéressées et aux autres entités}
\begin{itemize}
 \item La FCEG demande au BCAPG d'étudier la possibilité de modifier les critères d'accréditation afin de mieux tenir compte de la charge de travail des étudiants et d'exiger l'accès à des services de santé mentale de base pour tous les étudiants de premier cycle en ingénierie.
 \item La FCEG incite le BCAPG à étudier et à revoir les unités d'accréditation (UA) avec leurs heures respectives afin de s'assurer qu'elles restent réalisables dans le cadre d'une charge de travail complète pour les étudiants.
 \item La FCEG incite les DDIC à examiner la Norme nationale pour la santé mentale et le bien-être pour les étudiants du postsecondaire de la Commission de la santé mentale du Canada (Normes postsecondaires de la CSMC) [6] et à mettre en œuvre ces normes dans leurs facultés d'ingénierie.
 \item La FCEG incite le BCAPG à tenir ses instructeurs responsables de veiller à ce que la charge de travail ne nuise pas aux étudiants en ingénierie et à s'efforcer de mettre en place davantage de possibilités d'apprentissage tout au long de la vie, de possibilités d'apprentissage basées sur les problèmes et moins de défis académiques à fort enjeu.
 \item La FCEG incite Ingénieurs Canada et les DDIC de s'associer à la FCEG pour faire de la recherche liée à la santé mentale et à la charge de travail des étudiants dans le but de clarifier les impacts de la charge de travail sur la santé mentale des étudiants en ingénierie et les meilleures procédures pour l'améliorer.
 \item La FCEG incite ses organisations régionales partenaires (WESST, CAÉGO, CRÉIQ, CÉGA) à collaborer à des initiatives futures en matière de charge de travail et de santé mentale des étudiants, ce qui pourrait inclure un système de suivi de la charge de travail ou une deuxième enquête nationale auprès des étudiants.
 \item La FCEG incite ses membres à évaluer objectivement leur charge de travail et à trouver des lacunes ou des problèmes dans le cursus d'ingénierie de leurs écoles, et à soulever la question auprès de leurs organisations régionales.
\end{itemize}

\stanceheading{Ce que fait la FCEG}
\begin{itemize}
 \item La FCEG a mené une enquête nationale auprès des étudiants en 2023 afin de recueillir des données sur la santé mentale des étudiants en génie ainsi que sur la charge de travail et la perception des étudiants. La FCEG a organisé des séances et des conférences sur la santé mentale des étudiants, en mettant en évidence les initiatives étudiantes et l'épuisement des étudiants bénévoles, lors de ses principaux événements, notamment la Conférence canadienne sur le leadership en ingénierie (CCLI) et la Conférence sur la diversité en ingénierie (CDI).
 \item La FCEG a créé un répertoire des meilleures pratiques en matière de santé mentale disponible à toute organisation membre via le dossier \textit{Google Drive} externe.
\end{itemize}

\stanceheading{Ce que la FCEG envisage de faire}
\begin{itemize}
 \item La FCEG continuera à faciliter le partage des meilleures pratiques en matière de santé mentale entre ses sociétés membres et à mettre à jour les ressources pour la défense de la santé mentale des étudiants, le cas échéant.
 \item La FCEG poursuivra ses recherches sur les conséquences spécifiques d'une mauvaise santé mentale des étudiants et sur les meilleures pratiques pour améliorer la santé mentale dans les écoles d'ingénieurs.
 \item Le FCEG évaluera comment il peut mieux promouvoir la santé mentale des étudiants en interne par le biais des opérations et du contenu de ses événements.
 \item La FCEG assurera le suivi avec les parties prenantes concernées, y compris les DDIC, afin de garantir un examen systématique et régulier de la charge de travail et des horaires des programmes d'ingénierie, ainsi que des politiques mises en place pour protéger les étudiants.
 \item La FCEG examinera et appliquera les Normes Postsecondaires de la CSMC à la planification des conférences, à la motivation pour la défense des intérêts et aux pratiques de l'équipe interne.
\end{itemize}

\newpage
